\chapter{An Active-Set Method for the Regularized Master Problem}
\label{chap:qpas}
\chaptermark{Active-Set Method for the Master Problem}
\chaptermark{Active-Set Method for the Master Problem}

\vspace*{\fill}

This chapter sets up the \textbf{quadratic programming active-set method (QPAS)} for the
eliminated regularized master problem. We first derive its optimality conditions and show
that the Lagrange multipliers of the nonnegativity constraints are the reduced costs of the
columns and the duals of the inequality rows. The optimality test of the master problem and
the pricing test of column generation are therefore the same test.

\medskip

We then recall the primal active-set method: it fixes a \textbf{working set} of variables at
zero, minimizes over the remaining free variables, and adds or removes one constraint at a
time. Every iteration requires the solution of a KKT system whose constraint matrix consists
of the convexity rows and the rows of the working set.

\medskip

The main result of this chapter is a \textbf{reordering theorem} for this constraint matrix.
The columns of the regularized master problem are partitioned into five sets:
representatives, working-set columns, nullspace columns, active surpluses and free surpluses.
In this ordering the constraint matrix has a block structure with identity blocks, which
is the basis of the explicit nullspace method of the next chapter.

\vspace*{\fill}

\newpage


\section{The Master Problem in Simplified Notation}
\label{sec:simplified}

From now on we consider a single restricted master problem with $n$ columns. To keep the
notation light, we drop the explicit reference to the restricted index set and write
\begin{itemize}
    \item $x\in\R^n$ for the column weights (the vector $\lambda$ restricted to $\mathcal{J}'$),
    \item $A_E\in\R^{m_E\times n}$, $A_I\in\R^{m_I\times n}$, $A = [A_E;A_I]$ and
          $C\in\{0,1\}^{K\times n}$ for the corresponding columns of the linking and convexity
          matrices,
    \item $c\in\R^n$ for the costs of the columns, and $a_j\in\R^m$ for column $j$ of $A$,
    \item $y = [x; s]\in\R^{n+m_I}$ for the unified variable.
\end{itemize}
The vector $c$ always denotes the original column costs; the linear term of the quadratic
objective is $\ctil$ from~\eqref{eq:ctilde}. With this notation the eliminated regularized
master problem~\eqref{eq:eliminated} reads
\begin{subequations} \label{eq:qp}
\begin{align}
\min_{y\in\R^{n+m_I}}\quad & f(y) = \frac{1}{2\delta}\,y^\top\Abar^\top\Abar\,y + \ctil^\top y\\
\text{s.t.}\quad & \Cbar y = \one, \label{eq:qp-conv}\\
& y\ge0, \label{eq:qp-bounds}
\end{align}
\end{subequations}
with
\[
\Abar = \begin{bmatrix} A_E & 0\\ A_I & -I\end{bmatrix}\in\R^{m\times(n+m_I)},\qquad
\Cbar = \begin{bmatrix} C & 0\end{bmatrix}\in\R^{K\times(n+m_I)} .
\]
We write $g = \nabla f(y)$ for the gradient at the current iterate and, by
\cref{lem:gradient-dual},
\begin{equation} \label{eq:gradient-split}
g = \begin{bmatrix} g_x\\ g_s\end{bmatrix} = \begin{bmatrix} c - A^\top\pi\\ \pi_I\end{bmatrix},
\qquad \pi = \pi(y) = \pihat + \tfrac{1}{\delta}(b - \Abar y).
\end{equation}


\section{Optimality Conditions}
\label{sec:kkt}

Problem~\eqref{eq:qp} is a convex quadratic program with linear constraints, so the KKT
conditions are necessary and sufficient for optimality.

\begin{theorem}[KKT conditions of the regularized master problem] \label{thm:kkt-master}
A feasible point $y = [x;s]$ of~\eqref{eq:qp} is optimal if and only if there exist
$\mu\in\R^K$ and $\nu\in\R^{n+m_I}$ with
\[
g = \Cbar^\top\mu + \nu,\qquad \nu\ge0,\qquad \nu_\ell\,y_\ell = 0\quad\text{for all }\ell .
\]
Written per variable, with $\pi = \pi(y)$, these conditions are equivalent to:
\begin{enumerate}
    \item \textbf{reduced costs:} $\bar c_j(\pi,\mu) = c_j - a_j^\top\pi - \mu_{\psi(j)}\ge0$
          for every column $j$, with equality if $x_j>0$;
    \item \textbf{sign condition:} $\pi_i\ge0$ for every inequality row $i$;
    \item \textbf{slackness:} $\pi_i = 0$ for every inequality row $i$ with $s_i>0$.
\end{enumerate}
The multipliers are $\nu_j = \bar c_j(\pi,\mu)$ for the columns and $\nu_{n+i} = \pi_i$ for
the surpluses.
\end{theorem}

\begin{proof}
The first part is the KKT theorem for convex quadratic programs. For a column $j$, the
stationarity condition reads $g_j = \mu_{\psi(j)} + \nu_j$, so
$\nu_j = c_j - a_j^\top\pi - \mu_{\psi(j)} = \bar c_j(\pi,\mu)$ by~\eqref{eq:gradient-split}.
For surplus $i$ the column of $\Cbar$ is zero, so $\nu_{n+i} = g_{n+i} = \pi_i$. The
conditions $\nu\ge0$ and $\nu_\ell y_\ell = 0$ then become (i)--(iii). 
\end{proof}

\begin{remark}[The optimality test is the pricing test] \label{rem:same-test}
By \cref{thm:kkt-master}, the multipliers of the nonnegativity constraints of the columns
are their reduced costs, the same quantity that the pricing problem minimizes over the
columns that have not been generated yet. Solving the regularized master problem to
optimality establishes ``reduced cost $\ge0$'' for the columns inside the restricted problem,
and pricing establishes it for the columns outside. The multipliers of the surpluses carry
the sign conditions~\eqref{eq:dual-sign} of the inequality rows. Together with $\pi=\pihat$,
this is exactly the optimality certificate of \cref{thm:exact-termination}.
\end{remark}

\begin{remark}[Matching tolerances] \label{rem:tolerances}
The active-set method accepts a multiplier as nonnegative if $\nu_\ell\ge-\varepsilon$, and
pricing accepts a column only if $\bar c_j<-\varepsilon_{\mathrm{rc}}$. Choosing
$\varepsilon\ll\varepsilon_{\mathrm{rc}}$ (we use $\varepsilon = 10^{-10}\max(1,\max_j|c_j|)$
and $\varepsilon_{\mathrm{rc}} = 10^{-6}$) guarantees that pricing never returns a column that
is already in the restricted master problem.
\end{remark}


\section{The Primal Active-Set Method}
\label{sec:active-set}

\begin{definition}[Active set and working set] \label{def:working-set}
The \textbf{active set} of a feasible point $y$ is
$\mathcal{A}(y) = \{\ell : y_\ell = 0\}$. A \textbf{working set} is a set
$\Wset\subseteq\{1,\dots,n+m_I\}$ of indices that are treated as active, i.e.\ $y_\ell = 0$
for all $\ell\in\Wset$. We split it into the \textbf{column working set} $\Wx$ and the
\textbf{surplus working set} $\Ws$:
\[
\Wset = \Wx\cup\{\,n+i : i\in\Ws\,\},\qquad \Wx\subseteq\{1,\dots,n\},\quad
\Ws\subseteq\{1,\dots,m_I\} .
\]
A surplus $i\in\Ws$ is called \textbf{active}; a surplus $i\notin\Ws$ is called
\textbf{free}.
\end{definition}

Given a working set and a feasible iterate $y_k$ with $y_\ell = 0$ on $\Wset$, the
active-set method computes a search direction from the \textbf{equality-constrained
subproblem}
\begin{equation} \label{eq:eqp}
\min_{p}\;\frac{1}{2\delta}\,p^\top\Abar^\top\Abar\,p + g^\top p
\qquad\text{s.t.}\qquad \Cbar p = 0,\qquad p_\ell = 0\;\;(\ell\in\Wset),
\end{equation}
which minimizes $f(y_k+p)$ over the face of the feasible region defined by the working set.
To write its constraints compactly, let $W_x\in\R^{|\Wx|\times n}$ and
$W_s\in\R^{|\Ws|\times m_I}$ be the selection matrices whose rows are the unit vectors
$e_j^\top$, $j\in\Wx$, and $e_i^\top$, $i\in\Ws$.

\begin{definition}[Constraint matrix of a working set] \label{def:DW}
The \textbf{constraint matrix associated with $\Wset$} is
\begin{equation} \label{eq:DW}
D_\Wset = \begin{bmatrix} C & 0\\ W_x & 0\\ 0 & W_s\end{bmatrix}
\in\R^{(K+|\Wx|+|\Ws|)\times(n+m_I)} .
\end{equation}
The constraints of~\eqref{eq:eqp} are $D_\Wset\,p = 0$.
\end{definition}

The KKT conditions of~\eqref{eq:eqp} form the linear \textbf{KKT system}
\begin{equation} \label{eq:kkt-system}
\begin{bmatrix} \frac{1}{\delta}\Abar^\top\Abar & D_\Wset^\top\\ D_\Wset & 0\end{bmatrix}
\begin{bmatrix} p\\ -\mu\\ -\nu_x\\ -\nu_s\end{bmatrix}
= \begin{bmatrix} -g\\ 0\\ 0\\ 0\end{bmatrix},
\end{equation}
where $\mu\in\R^K$, $\nu_x\in\R^{|\Wx|}$ and $\nu_s\in\R^{|\Ws|}$ are the Lagrange
multipliers of the convexity constraints and of the working-set constraints of the columns
and of the surpluses. If $p = 0$ solves~\eqref{eq:kkt-system}, the iterate minimizes $f$
over its face, and we call it a \textbf{subspace minimizer}. The generic method is
\cref{alg:qpas}.

\begin{algobox}{Primal Active-Set Method for the Regularized Master Problem}{alg:qpas}
\State \textbf{Input:} feasible $y$, working set $\Wset$ with $y_\ell = 0$ on $\Wset$
\Repeat
    \State Solve the KKT system~\eqref{eq:kkt-system} for $p$, $\mu$, $\nu = (\nu_x,\nu_s)$
    \If{$p = 0$}
        \If{$\nu_\ell\ge-\varepsilon$ for all $\ell\in\Wset$}
            \State \textbf{stop}: $y$ is optimal
        \Else
            \State $r\gets\arg\min_{\ell\in\Wset}\nu_\ell$;\quad $\Wset\gets\Wset\setminus\{r\}$
                   \Comment{reduce the working set}
        \EndIf
    \Else
        \State $\alpha\gets\min\Bigl\{1,\;\min\limits_{\ell\notin\Wset,\;p_\ell<0}\;
               -\dfrac{y_\ell}{p_\ell}\Bigr\}$ \Comment{ratio test}
        \State $y\gets y+\alpha p$
        \If{$\alpha<1$}
            \State $q\gets$ the index attaining the minimum;\quad $\Wset\gets\Wset\cup\{q\}$
                   \Comment{expand the working set}
        \EndIf
    \EndIf
\Until{stopped}
\end{algobox}

Under the usual nondegeneracy assumptions, every iteration either decreases $f$ strictly or
changes the working set in a way that cannot cycle, so the method terminates after finitely
many iterations \cite{nocedal2006numerical, fletcher1971quadratic}. Its cost is dominated by
the solution of the KKT system~\eqref{eq:kkt-system}, which has dimension
$n + m_I + K + |\Wset|$ and grows with every generated column. The rest of this work shows
how to avoid ever forming it. The nullspace method of \cref{chap:nullspace} needs two
properties:

\begin{assumption} \label{ass:nullspace}
\begin{enumerate}
    \item $D_\Wset$ has full row rank;
    \item the reduced Hessian $Z^\top\Abar^\top\Abar Z$ is positive definite, where the columns
          of $Z$ form a basis of $\Null(D_\Wset)$.
\end{enumerate}
\end{assumption}

Property~(i) holds for every working set that keeps one column per block free
(\cref{thm:rank} below). Property~(ii) can fail, and \cref{chap:singular} shows how the
method guarantees it.


\section{The Constraint Matrix and the Reordering Theorem}
\label{sec:reordering}

The structure of $D_\Wset$ is determined by the convexity matrix $C$, which has the column
generation property, and by the zero columns of the surpluses, which do not. We first
arrange the columns of the weights $x$, and then add the surpluses.

\begin{definition}[Free columns] \label{def:free-columns}
The \textbf{free columns} are $\Fset_x = \{1,\dots,n\}\setminus\Wx$.
\end{definition}

\begin{lemma} \label{lem:free-block}
If $y$ is feasible and $x_j = 0$ for all $j\in\Wx$, then every block has at least one free
column, and the submatrix $C_{\Fset_x}$ of $C$ has the column generation property.
\end{lemma}

\begin{proof}
For every block $k$, $\sum_{\psi(j) = k}x_j = 1$ and $x_j = 0$ on $\Wx$, so some column of
block $k$ lies in $\Fset_x$. The claim follows from \cref{lem:cg-property}. 
\end{proof}

\begin{definition}[Representatives] \label{def:representatives}
A \textbf{list of representatives} is an ordered list $\Sset = [S_1,\dots,S_K]$ of free
columns with $\psi(S_k) = k$, one per block. The \textbf{representative of column} $j$ is
\[
\rep(j) = S_{\psi(j)} .
\]
The remaining free columns form the \textbf{nullspace set} $\Nx = \Fset_x\setminus\Sset$.
\end{definition}

The representatives exist by \cref{lem:free-block}. Their columns in $C$ are the unit vectors
$e_1,\dots,e_K$, so $C_\Sset = I_K$.

\begin{definition}[Extended ordering] \label{def:extended-ordering}
The \textbf{extended ordering} of the variables of~\eqref{eq:qp} is the partition into the
five sets
\begin{center}
\renewcommand{\arraystretch}{1.25}
\begin{tabular}{lll}
\toprule
set & contents & status \\
\midrule
$\Sset$ & one representative column per block & free \\
$\Wx$   & columns in the working set & $x_j = 0$ fixed \\
$\Nx$   & the remaining free columns & free \\
$\Ws$   & active surpluses & $s_i = 0$ fixed \\
$\Ns$   & free surpluses & free \\
\bottomrule
\end{tabular}
\end{center}
We write $[\,\Sset,\Wx,\Nx\mid\Ws,\Ns\,]$ for this ordering.
\end{definition}

\begin{theorem}[Reordering theorem] \label{thm:reordering}
Let $\Wset = \Wx\cup\Ws$ be a working set such that every block has a free column, and let
$[\,\Sset,\Wx,\Nx\mid\Ws,\Ns\,]$ be an extended ordering. Let $Q$ be the permutation matrix
that orders the variables in this way. Then
\begin{equation} \label{eq:reordered}
D_\Wset\,Q =
\begin{bmatrix}
I_K & C_{\Wx} & C_{\Nx} & 0 & 0\\
0 & I_{|\Wx|} & 0 & 0 & 0\\
0 & 0 & 0 & I_{|\Ws|} & 0
\end{bmatrix},
\end{equation}
where $C_{\Wx}$ and $C_{\Nx}$ are the columns of $C$ indexed by $\Wx$ and $\Nx$.
\end{theorem}

\begin{proof}
The permutation is block-diagonal, $Q = \diag(Q_x, Q_s)$: $Q_x$ orders the column weights as
$[\Sset,\Wx,\Nx]$ and $Q_s$ orders the surpluses as $[\Ws,\Ns]$. Apply it to the three block
rows of~\eqref{eq:DW}.
\begin{itemize}
    \item The convexity rows $[\,C\;\;0\,]$ become $[\,C_\Sset\;\;C_{\Wx}\;\;C_{\Nx}\;\;0\;\;0\,]$,
          and $C_\Sset = I_K$ by \cref{def:representatives}.
    \item The rows $[\,W_x\;\;0\,]$ select the columns in $\Wx$, which appear in the second
          position: $[\,0\;\;I_{|\Wx|}\;\;0\;\;0\;\;0\,]$.
    \item The rows $[\,0\;\;W_s\,]$ select the surpluses in $\Ws$, which appear in the fourth
          position: $[\,0\;\;0\;\;0\;\;I_{|\Ws|}\;\;0\,]$. 
\end{itemize}
\end{proof}

\begin{figure}[ht]
\centering
\begin{tikzpicture}[x=0.6cm, y=0.6cm, font=\small]
    % column group widths: S=2, Wx=6, Nx=3, Ws=2, Ns=2 ; rows: K=2, Wx=6, Ws=2
    \def\Sx{0}\def\Wxx{2}\def\Nxx{8}\def\Wsx{11}\def\Nsx{13}\def\Ed{15}
    \def\Tp{10}\def\rK{8}\def\rW{2}\def\rB{0}
    % shaded nonzero blocks
    \fill[cblue!35]   (\Sx,\rK)  rectangle (\Wxx,\Tp);
    \fill[cblue!12]   (\Wxx,\rK) rectangle (\Nxx,\Tp);
    \fill[cblue!12]   (\Nxx,\rK) rectangle (\Wsx,\Tp);
    \fill[corange!35] (\Wxx,\rW) rectangle (\Nxx,\rK);
    \fill[cgreen!35]  (\Wsx,\rB) rectangle (\Nsx,\rW);
    % identity diagonals
    \draw[very thick, cblue!80!black]   (\Sx,\Tp) -- (\Wxx,\rK);
    \draw[very thick, corange!80!black] (\Wxx,\rK) -- (\Nxx,\rW);
    \draw[very thick, cgreen!60!black]  (\Wsx,\rW) -- (\Nsx,\rB);
    % grid
    \draw[thick] (\Sx,\rB) rectangle (\Ed,\Tp);
    \foreach \c in {\Wxx,\Nxx,\Wsx,\Nsx} \draw (\c,\rB) -- (\c,\Tp);
    \draw[thick, dashed] (\Wsx,\rB) -- (\Wsx,\Tp);
    \foreach \r in {\rK,\rW} \draw (\Sx,\r) -- (\Ed,\r);
    % labels inside
    \node at (1.4,9.4) {\scriptsize $I_K$};
    \node at (5,9) {$C_{\Wx}$};
    \node at (9.5,9) {$C_{\Nx}$};
    \node at (12,9) {$0$}; \node at (14,9) {$0$};
    \node[fill=white, inner sep=1pt] at (5.6,4.6) {$I_{|\Wx|}$};
    \node at (12,1) [right=-2pt] {};
    \node[fill=white, inner sep=1pt] at (12.3,0.6) {\scriptsize $I_{|\Ws|}$};
    % column headers
    \node at (1,10.6) {$\Sset$};
    \node at (5,10.6) {$\Wx$};
    \node at (9.5,10.6) {$\Nx$};
    \node at (12,10.6) {$\Ws$};
    \node at (14,10.6) {$\Ns$};
    \draw[decorate, decoration={brace, amplitude=4pt}] (0,11.2) -- (11,11.2)
        node[midway, above=4pt] {column weights $x$};
    \draw[decorate, decoration={brace, amplitude=4pt}] (11,11.2) -- (15,11.2)
        node[midway, above=4pt] {surpluses $s$};
    % row labels
    \node[left] at (0,9) {$C$};
    \node[left] at (0,5) {$W_x$};
    \node[left] at (0,1) {$W_s$};
\end{tikzpicture}
\caption{Block structure of the reordered constraint matrix $D_\Wset Q$ of
\cref{thm:reordering}. The diagonal lines are identity blocks. The surplus part (right of the
dashed line) only contributes the identity block of the active surpluses; the free
surpluses $\Ns$ have empty columns.}
\label{fig:reordered}
\end{figure}

\begin{theorem}[Full row rank] \label{thm:rank}
Under the assumptions of \cref{thm:reordering}, $D_\Wset$ has full row rank,
\[
\rank(D_\Wset) = K + |\Wx| + |\Ws|,
\]
and $\dim\Null(D_\Wset) = |\Nx| + |\Ns|$.
\end{theorem}

\begin{proof}
In~\eqref{eq:reordered}, the identity blocks $I_K$, $I_{|\Wx|}$ and $I_{|\Ws|}$ lie in
disjoint rows and disjoint columns, so the rows are linearly independent. Permutations do not
change the rank. The dimension of the nullspace follows from the rank--nullity theorem:
$(n+m_I) - (K+|\Wx|+|\Ws|) = |\Nx| + |\Ns|$, because $n = K + |\Wx| + |\Nx|$ and
$m_I = |\Ws| + |\Ns|$. 
\end{proof}

\Cref{thm:rank} establishes property~(i) of \cref{ass:nullspace} for every working set the
method can produce: the ratio test never removes the last free column of a block, since the
representative of the block can only reach zero while some other column of the block grows
(\cref{lem:replacement} below).

\begin{remark}[Comparison with problems without inequality rows] \label{rem:three-sets}
Without inequality rows ($m_I = 0$), the surplus sets are empty, and the extended ordering
reduces to the three sets $[\Sset,\Wx,\Nx]$. The inequality rows thus only add two groups of
columns: active surpluses, which behave like working-set columns, and free surpluses, which
behave like nullspace columns without a representative. The loss of the column generation
property (\cref{rem:cg-property-lost}) does not affect the reordering. It only means that the
surpluses form groups of their own.
\end{remark}

\Cref{fig:spy-DW} shows the effect of the reordering on an actual constraint matrix, taken
from the final working set of a column generation run on one of the test instances of
\cref{chap:results}.

\begin{figure}[ht]
\centering
\IfFileExists{data/spy_DW_orig.csv}{%
\begin{tikzpicture}
\begin{groupplot}[group style={group size=2 by 1, horizontal sep=1.2cm},
    width=0.5\textwidth, height=0.42\textwidth, y dir=reverse,
    xmin=0, ymin=0, enlargelimits=false, tick label style={font=\scriptsize},
    title style={font=\small}, xlabel style={font=\small}, ylabel style={font=\small}]
\nextgroupplot[title={original order}, xlabel={variable}, ylabel={row of $D_\Wset$}]
    \addplot[only marks, mark=*, mark size=0.25pt, cblue]
        table[x=col, y=row, col sep=comma] {data/spy_DW_orig.csv};
\nextgroupplot[title={order $[\Sset,\Wx,\Nx\mid\Ws,\Ns]$}, xlabel={variable}]
    \addplot[only marks, mark=*, mark size=0.25pt, cblue]
        table[x=col, y=row, col sep=comma] {data/spy_DW_perm.csv};
\end{groupplot}
\end{tikzpicture}}{\fbox{spy plot data missing: run scripts/make\_plot\_data.jl}}
\caption{Sparsity pattern of the constraint matrix $D_\Wset$ at the end of a column
generation run (instance I2 with 30\% inequality rows), before and after the reordering of
\cref{thm:reordering}. After reordering, the identity blocks of the representatives (top
left), the working-set columns (diagonal) and the active surpluses (bottom right) are
visible.}
\label{fig:spy-DW}
\end{figure}
